\section{Free Haar unitaries in an ultrapower}
\label{sec:independence}

Throughout this section, $P\subset M$ is an irreducible inclusion of
$\mathrm{II}_1$ factors.  Fix a free ultrafilter $\omega$ on
$\mathbb N$.  We use the tracial
ultrapower
\[
  \mathbf M=M^\omega
  =\ell^\infty(\mathbb N,M)/\mathcal I_\omega,
  \qquad
  \mathcal I_\omega
  =\left\{(b_n)\in\ell^\infty(\mathbb N,M):
       \lim_{n\to\omega}\|b_n\|_2=0\right\},
\]
and its von Neumann subalgebra $\mathbf Q=P^\omega$.  The trace of an
element represented by $(b_n)$ is $\lim_{n\to\omega}\tau(b_n)$; we
denote this trace, and its associated $2$-norm, by the same symbols as
on $M$.  We identify $M$ with its constant copy in $\mathbf M$.

A unitary $t$ in a tracial algebra is called \emph{Haar} if
$\tau(t^n)=0$ for every nonzero integer $n$.  Two unital subalgebras
are \emph{free} if every alternating product of centered elements
from the two subalgebras has trace zero.  We say that a unitary is
free from a unital algebra $D$ when $C^*(t)$ and $D$ are free.

For the perturbation in Section~\ref{sec:perturbation}, we will choose
successive Haar unitaries in $\mathbf Q$, each free from the algebra
containing the prescribed data and the earlier choices.  We will also
need this freeness to survive multiplication on either side by unitaries
from that algebra.  Popa's theorem supplies independence relative to
$\mathbf Q'\cap\mathbf M$, so we first identify that commutant.

\begin{lemma}
\label{lem:ultrapower-commutant}
The inclusion $\mathbf Q\subset\mathbf M$ is irreducible:
\[
  \mathbf Q'\cap\mathbf M=\mathbb C1.
\]
\end{lemma}

\begin{proof}
Let $b\in\mathbf Q'\cap\mathbf M$ have a uniformly bounded
representative $(b_n)$, and put
\[
  \delta_n=\sup_{w\in\mathcal U(P)}
       \|wb_nw^*-b_n\|_2.
\]
Then $\lim_{n\to\omega}\delta_n=0$.  Otherwise there would be an
$\varepsilon>0$ and a set belonging to $\omega$ on which
$\delta_n>\varepsilon$.  Choosing on this set a unitary
$w_n\in P$ with $\|w_nb_nw_n^*-b_n\|_2>\varepsilon$, and setting
$w_n=1$ elsewhere, would give a unitary in $\mathbf Q$ that does not
commute with $b$.

For each $n$, let $K_n$ be the closed convex hull, in $L^2(M,\tau)$,
of the elements $wb_nw^*$, $w\in\mathcal U(P)$.  The unique element
$c_n\in K_n$ of smallest $2$-norm is fixed by conjugation by every
unitary of $P$: each such conjugation preserves $K_n$ and its Hilbert
space norm.  Moreover, $c_n$ belongs to $M$ and
$\|c_n\|\leq\|b_n\|$.  Indeed, convex combinations from this orbit
have this operator norm bound, and an operator-norm-bounded
$2$-norm limit lies in the same bounded ball of $M$.  To see the latter
fact, left multiplication by such a convergent sequence converges
strongly on the dense subspace $M\subset L^2(M)$, since
\[
  \|(x_r-x_s)a\|_2\leq\|x_r-x_s\|_2\|a\|
  \qquad(a\in M),
\]
and hence converges strongly on all of $L^2(M)$ by the uniform
operator norm bound.  The strong limit belongs to $M$, and its value
at $1$ is the prescribed $2$-norm limit.

Thus $c_n\in P'\cap M=\mathbb C1$.  The trace is constant on $K_n$
and continuous in $2$-norm, so $c_n=\tau(b_n)1$.  Every element of
$K_n$ is at $2$-distance at most $\delta_n$ from $b_n$.  Consequently
\[
  \|b_n-\tau(b_n)1\|_2\leq\delta_n\longrightarrow_{n\to\omega}0.
\]
The bounded scalar sequence $(\tau(b_n))$ has a scalar ultralimit,
which represents $b$.  This proves the assertion.
\end{proof}

We use the following special case of Popa's independence theorem.
The atomic-relative-commutant case needed here already appears in
\cite[Lemma~1.4 and Theorem~2.1]{Popa1995}; we state the later
formulation from~\cite[Theorem~0.1(a)]{Popa2014}.

\begin{theorem}[Popa's independence theorem]
\label{thm:popa-independence}
Let $P\subset M$ be an irreducible inclusion of $\mathrm{II}_1$
factors, and put $\mathbf Q=P^\omega\subset M^\omega=\mathbf M$.
For every linear subspace $X\subset\mathbf M$ that is separable in
$2$-norm and satisfies $\tau(x)=0$ for all $x\in X$, there is a
diffuse abelian von Neumann algebra $A\subset\mathbf Q$ such that
\begin{equation}
\label{eq:popa-independence}
  \tau(x_0a_1x_1\cdots a_mx_m)=0
\end{equation}
whenever $m\geq1$, $x_0\in X\cup\{1\}$,
$x_1,\ldots,x_m\in X$, and $a_1,\ldots,a_m\in A$ have trace zero.
\end{theorem}

This is the constant-inclusion case of
\cite[Theorem~0.1(a)]{Popa2014}.  The hypothesis there is that no
nonzero corner of the coordinate algebra $P$ intertwines into
$P'\cap M$ inside $M$.  It holds here because $P'\cap M=\mathbb C1$:
a nonzero corner of the diffuse factor $P$ cannot embed normally
into a finite matrix amplification of the scalars.  The conclusion
of that theorem is independence relative to
$\mathbf Q'\cap\mathbf M$.  Lemma~\ref{lem:ultrapower-commutant}
identifies this commutant with the scalars, giving
\eqref{eq:popa-independence}.

\begin{lemma}
\label{lem:free-haar-extension}
Let $D\subset\mathbf M$ be a norm-separable unital $C^*$-algebra.
There is a Haar unitary $t\in\mathbf Q$ free from $D$.  In
particular,
\begin{equation}
\label{eq:free-endpoints}
  \tau(d_0t^{n_1}d_1\cdots t^{n_j}d_j)=0
\end{equation}
for every $j\geq1$, every $n_1,\ldots,n_j\in\mathbb Z\setminus\{0\}$,
and all $d_0,\ldots,d_j\in D$ such that
$\tau(d_i)=0$ for $1\leq i\leq j-1$.  The two endpoints
$d_0,d_j$ are arbitrary.
\end{lemma}

\begin{proof}
The subspace $X=D\cap\ker\tau$ is norm separable, hence $2$-norm
separable.  For example, applying $d\mapsto d-\tau(d)1$ to a
countable norm-dense subset of $D$ gives a countable norm-dense
subset of $X$.  Apply Theorem~\ref{thm:popa-independence} to obtain
a diffuse abelian algebra $A\subset\mathbf Q$.

We first choose a Haar unitary in $A$.  Diffuseness gives arbitrarily
small nonzero subprojections of every nonzero projection: repeatedly
split and retain a piece of at most half the trace.  Normality of the
trace and maximality under inclusion then give subprojections of every
prescribed smaller trace.  We may therefore choose recursively
partitions of $1$ by projections
\[
  \{p_{m,k}:0\leq k<2^m\}\subset A,
  \qquad \tau(p_{m,k})=2^{-m},
\]
with $p_{m,k}=p_{m+1,2k}+p_{m+1,2k+1}$.  The self-adjoint
elements
\[
  h_m=\sum_{k=0}^{2^m-1}\frac{k}{2^m}p_{m,k}
\]
satisfy $\|h_{m+1}-h_m\|\leq2^{-(m+1)}$ and converge in norm to
an element $h\in A$.  For each continuous function $f$ on $[0,1]$,
\[
  \tau(f(h))
  =\lim_{m\to\infty}2^{-m}
        \sum_{k=0}^{2^m-1}f(k2^{-m})
  =\int_0^1 f(s)\,ds.
\]
It follows that $t=\exp(2\pi i h)$ is Haar.

We next remove the centering requirement at the two endpoints, so
that \eqref{eq:free-endpoints} can be applied to words with arbitrary
outside coefficients.  Expand each endpoint into its scalar and
centered parts.  A term with a centered last endpoint
vanishes by \eqref{eq:popa-independence}, since every nonzero power
of $t$ belongs to $A$ and is centered.  A term with a centered first
endpoint and scalar last endpoint vanishes after moving the first
endpoint to the end by traciality.  It remains to consider the word
with both endpoints omitted.  When $j=1$ this is the Haar property.
When $j\geq2$, rotate the final power $t^{n_j}$ next to the first
power $t^{n_1}$.  If $n_j+n_1\ne0$, the resulting word is covered
by \eqref{eq:popa-independence}.  If $n_j+n_1=0$, the two powers
cancel and leave a word starting and ending with centered elements
of $D$.  This word is again covered by
\eqref{eq:popa-independence}, except when $j=2$, in which case it
is simply the centered element $d_1$.  This proves
\eqref{eq:free-endpoints}.

Laurent polynomials in $t$ are norm dense in $C^*(t)$.  As $t$ is
Haar, the centered Laurent polynomials are precisely the linear
combinations of its nonzero powers.  Expanding an alternating
centered word, and using \eqref{eq:free-endpoints} followed by norm
approximation, proves that $t$ is free from $D$.
\end{proof}

\begin{lemma}[Absorption of base-algebra unitaries]
\label{lem:haar-absorption}
Let $D$ be a unital $C^*$-subalgebra of a finite tracial von Neumann
algebra, and let $t$ be a Haar unitary free from $D$.  If
$a,b\in\mathcal U(D)$, then both $atb$ and $t^*$ are Haar unitaries
free from $D$.  Both satisfy \eqref{eq:free-endpoints} in place of $t$.
\end{lemma}

\begin{proof}
For any Haar unitary free from $D$, formula
\eqref{eq:free-endpoints} follows immediately by expanding the
endpoints into scalar and centered parts and applying freeness to
each resulting alternating word.  Replacing $t$ by $t^*$ merely
negates the exponents, so we concentrate on $v=atb$.

Consider a word
\[
  d_0v^{n_1}d_1\cdots v^{n_j}d_j
\]
with nonzero integer exponents and centered internal factors
$d_1,\ldots,d_{j-1}$.  Expand each power of $v$ into individual
letters $t$ and $t^{-1}$, with elements of $D$ between them.  Within
one power of $v$ all these letters have the same sign.  A gap
between opposite signs can therefore occur only between two
successive powers of $v$.  Such a gap equals
\[
  b d_i b^* \quad\text{for a transition from $t$ to $t^{-1}$},
  \qquad
  a^*d_i a \quad\text{for a transition from $t^{-1}$ to $t$}.
\]
It is centered by traciality.  Every other gap joins equal signs.
Expand each of these latter gaps into its scalar and centered
parts.  In a resulting term, removing a scalar gap only joins
letters of equal sign.  Thus each block joined across scalar gaps
is a nonzero power of $t$, and consecutive blocks are separated by
centered elements of $D$.  The endpoint factors remain arbitrary
elements of $D$.  Formula~\eqref{eq:free-endpoints} shows that every
resulting term has trace zero.

We have proved \eqref{eq:free-endpoints} for $v$.  Taking $j=1$
and $d_0=d_1=1$ proves that $v$ is Haar.  The general formula and
norm approximation by Laurent polynomials then prove its freeness
from $D$.
\end{proof}

We may apply Lemma~\ref{lem:free-haar-extension} successively.
Specifically, given a norm-separable unital $C^*$-algebra $D_0\subset
\mathbf M$ and an integer $r\geq0$, choose $t_j\in\mathbf Q$ Haar
and free from $D_j$, and set
\[
  D_{j+1}=C^*(D_j,t_j),\qquad 0\leq j\leq r.
\]
Each $D_{j+1}$ remains norm separable: polynomials in a countable
norm-dense subset of $D_j$, the two elements $t_j,t_j^*$, and
rational complex coefficients form a countable dense set.  Thus
every application of the lemma has the required separability
hypothesis, even though the previously chosen unitaries already
belong to $\mathbf Q$.
